\documentclass[UTF8]{ctexart}
\usepackage{listings}
\usepackage{xcolor}


\usepackage{geometry}
\usepackage{listings}
\usepackage{xcolor}
\lstset{
    basicstyle=\ttfamily\small,
    breaklines=true,
    frame=single,
    columns=fullflexible,
    extendedchars=false,
    inputencoding=utf8,
    keepspaces=true
}

\geometry{a4paper, margin=2cm}

\lstset{
basicstyle=\ttfamily\small,
breaklines=true,
frame=single,
columns=fullflexible
}

\title{基于 WSL + rsync 的跨主机大文件传输完整指南}
\author{}
\date{}

\begin{document}

\maketitle

\section{总体架构}

本方案采用如下结构：

\begin{lstlisting}
发送端 WSL (rsync)
↓ SSH
接收端 Windows (OpenSSH + cwRsync)
\end{lstlisting}

特点：
\begin{itemize}
\item 不需要端口转发
\item 支持断点续传
\item 适用于大文件
\end{itemize}

---

\section{接收端配置（Windows）}

\subsection{1. 安装 cwRsync}

将下载的压缩包解压到：

\begin{lstlisting}
D:\cwrsync_6.4.7_x64_free
\end{lstlisting}

确认存在：

\begin{lstlisting}
D:\cwrsync_6.4.7_x64_free\bin\rsync.exe
\end{lstlisting}

---

\subsection{2. 验证 rsync（在 PowerShell 执行）}

\begin{lstlisting}
& "D:\cwrsync_6.4.7_x64_free\bin\rsync.exe" --version
\end{lstlisting}

若输出：

\begin{lstlisting}
rsync  version 3.x.x
\end{lstlisting}

说明安装成功。

---

\subsection{3. 启动 SSH 服务（PowerShell 执行）}

\begin{lstlisting}
Start-Service sshd
\end{lstlisting}

验证：

\begin{lstlisting}
Get-Service sshd
\end{lstlisting}

状态应为：

\begin{lstlisting}
Running
\end{lstlisting}

---

\subsection{4. 创建目标目录}

\begin{lstlisting}
mkdir D:\targetprofile
\end{lstlisting}

---

\section{发送端配置（WSL）}

\subsection{1. 安装 rsync}

在 WSL 终端执行：

\begin{lstlisting}
sudo apt update
sudo apt install rsync -y
\end{lstlisting}

---

\subsection{2. 验证 rsync}

\begin{lstlisting}
rsync --version
\end{lstlisting}

---

\subsection{3. 测试 SSH 连接}

\begin{lstlisting}
ssh ggbond@100.102.3.10
\end{lstlisting}

成功标志：

\begin{itemize}
\item 能进入远程终端
\end{itemize}

退出：

\begin{lstlisting}
exit
\end{lstlisting}

---

\section{文件路径说明}

\begin{itemize}
\item WSL 路径：
\end{itemize}

\begin{lstlisting}
/mnt/f/xxx
\end{lstlisting}

\begin{itemize}
\item Windows rsync 路径：
\end{itemize}

\begin{lstlisting}
/cygdrive/d/targetprofile/
\end{lstlisting}

---

\section{执行 rsync（发送端 WSL 执行）}

\subsection{完整命令}

\begin{lstlisting}
rsync -avP --inplace --partial 
-e "ssh -T -c [aes128-gcm@openssh.com](mailto:aes128-gcm@openssh.com) -o Compression=no" 
--rsync-path="D:/cwrsync_6.4.7_x64_free/bin/rsync.exe" 
"/mnt/f/[Mumuy] 金瓶梅 1-28(武松篇章完)+番外更新 (圣诞礼物 & 简中版 DLsite发售中！) [Chinese] [Ongoing]/" 
ggbond@100.102.3.10:"/cygdrive/d/漫画/"
\end{lstlisting}

---

\subsection{执行说明}

\begin{itemize}
\item 在 \textbf{发送端 WSL} 执行
\item 不需要先 ssh 登录
\end{itemize}

---

\subsection{首次执行提示}

\begin{lstlisting}
Are you sure you want to continue connecting (yes/no)?
\end{lstlisting}

输入：

\begin{lstlisting}
yes
\end{lstlisting}

然后输入密码（不会显示字符）

---

\section{运行结果}

成功时输出：

\begin{lstlisting}
sending incremental file list
xxx.zip
30%  20MB/s
\end{lstlisting}

---

\section{断点续传}

重新执行相同命令即可：

\begin{lstlisting}
rsync -avP ...
\end{lstlisting}

---
\section{日常使用指南}

本节用于说明系统搭建完成后，日常仅进行文件传输时，发送端与接收端分别需要开启或关闭哪些服务，以及推荐的日常操作流程。

\subsection{总体原则}

本方案的传输结构为：

\begin{lstlisting}
发送端 WSL (rsync)
      ↓ SSH
接收端 Windows (OpenSSH + cwRsync)
\end{lstlisting}

日常使用时应遵循以下原则：

\begin{itemize}
    \item 接收端必须开启 SSH 服务；
    \item 发送端无需开启 SSH 服务；
    \item 发送端只需打开 WSL 并执行 rsync 命令；
    \item 文件传输结束后，可根据安全需要关闭接收端 SSH 服务。
\end{itemize}

\subsection{开机后接收端需要执行的操作}

接收端是指接收文件的 Windows 电脑。

\subsubsection{2.1 启动 SSH 服务}

在接收端 \textbf{Windows PowerShell} 中执行：

\begin{lstlisting}
Start-Service sshd
\end{lstlisting}

\subsubsection{检查 SSH 服务是否正常运行}

在接收端 \textbf{Windows PowerShell} 中执行：

\begin{lstlisting}
Get-Service sshd
\end{lstlisting}

若输出中包含如下状态，则说明 SSH 服务已成功启动：

\begin{lstlisting}
Running
\end{lstlisting}

\subsubsection{ 检查 cwRsync 是否可用（可选）}

在接收端 \textbf{Windows PowerShell} 中执行：

\begin{lstlisting}
& "D:\cwrsync_6.4.7_x64_free\bin\rsync.exe" --version
\end{lstlisting}

若输出中包含 rsync 版本号，则说明接收端 rsync 可正常使用。

\subsubsection{ 确认 Tailscale 已连接（若通过 Tailscale 传输）}

若本方案通过 Tailscale 进行跨设备连接，则接收端还应确认 Tailscale 客户端已处于在线状态，并确保其 Tailscale IP 未发生异常变化。

\subsection{开机后发送端需要执行的操作}

发送端是指发起文件传输的电脑。

\subsubsection{打开 WSL 终端}

发送端无需启动 SSH 服务，只需打开 WSL 终端即可。

\subsubsection{检查 rsync 是否可用}

在发送端 \textbf{WSL 终端} 中执行：

\begin{lstlisting}
rsync --version
\end{lstlisting}

若输出中包含 rsync 版本号，则说明发送端 rsync 可正常使用。

\subsubsection{测试与接收端的 SSH 连通性（可选）}

在发送端 \textbf{WSL 终端} 中执行（这里的ip和用户名要根据tailscale里的写）：

\begin{lstlisting}
ssh ggbond@100.102.3.10
\end{lstlisting}

若能正常进入远程终端，则说明网络与 SSH 服务均正常。测试完成后可执行：

\begin{lstlisting}
exit
\end{lstlisting}

退出远程终端。

\subsection{日常传输文件的操作步骤}

\subsubsection{在发送端 WSL 中执行 rsync 命令}

在发送端 \textbf{WSL 终端} 中直接执行下列命令：

\begin{lstlisting}
rsync -avP --inplace --partial \
-e "ssh -T -c aes128-gcm@openssh.com -o Compression=no" \
--rsync-path="D:/cwrsync_6.4.7_x64_free/bin/rsync.exe" \
"/mnt/f/yourprofile/" \
ggbond@100.102.3.10:"/cygdrive/d/targetprofile/"
\end{lstlisting}

需要注意：

\begin{itemize}
    \item 该命令必须在发送端 \textbf{本地 WSL} 中执行；
    \item 不需要先手动执行 \texttt{ssh ggbond@100.102.3.10} 再执行 rsync；
    \item rsync 会自动调用 SSH 完成远程连接。
\end{itemize}

\subsubsection{首次连接时的确认}

若为第一次连接接收端，系统可能提示：

\begin{lstlisting}
Are you sure you want to continue connecting (yes/no)?
\end{lstlisting}

此时输入：

\begin{lstlisting}
yes
\end{lstlisting}

然后回车。

\subsubsection{输入密码}

随后系统会提示输入密码。输入时终端中不会显示任何字符，此为正常现象。输入完成后直接回车即可。

\subsubsection{传输中断后的续传}

若传输过程中中断，只需在发送端 WSL 中重新执行同一条 rsync 命令即可。由于已启用 \texttt{-P} 与 \texttt{--partial} 参数，rsync 会自动从中断位置继续传输。

\subsection{文件传输完成后需要关闭哪些服务}

\subsubsection{发送端}

发送端在文件传输完成后，无需执行额外关闭操作。直接关闭 WSL 终端即可。

\subsubsection{接收端}

若出于安全考虑，不希望接收端长期开放 SSH 服务，则可在接收端 \textbf{Windows PowerShell} 中执行：

\begin{lstlisting}
Stop-Service sshd
\end{lstlisting}

这样可在不使用时关闭 SSH 服务。

\subsubsection{下次使用时重新开启}

下次需要传输文件时，仅需再次在接收端执行：

\begin{lstlisting}
Start-Service sshd
\end{lstlisting}

即可恢复使用。

\subsection{推荐的日常最简流程}

日常使用时，建议采用如下最简操作流程：

\begin{enumerate}
    \item 接收端开机后，在 Windows PowerShell 中执行：
\begin{lstlisting}
Start-Service sshd
\end{lstlisting}

    \item 发送端打开 WSL，直接执行 rsync 命令传输文件；
\begin{lstlisting}
rsync -avP --inplace --partial \
-e "ssh -T -c aes128-gcm@openssh.com -o Compression=no" \
--rsync-path="D:/cwrsync_6.4.7_x64_free/bin/rsync.exe" \
"/mnt/f/yourprofile/" \
ggbond@100.102.3.10:"/cygdrive/d/targetprofile/"
\end{lstlisting}
    \item 文件传输完成后，如需提高安全性，则在接收端执行：
\begin{lstlisting}
Stop-Service sshd
\end{lstlisting}

    \item 查看sshd状态：
\begin{lstlisting}
Get-Service sshd
\end{lstlisting}
\end{enumerate} 
\subsection{日常使用要点总结}

\begin{itemize}
    \item 接收端必须开启 SSH 服务；
    \item 发送端不需要开启 SSH 服务；
    \item rsync 命令始终在发送端 WSL 中执行；
    \item 传输结束后，接收端 SSH 服务可按需关闭；
    \item 再次传输时，仅需重新开启接收端 SSH 服务。
\end{itemize}

\section{总结}

\begin{itemize}
\item 发送端：WSL + rsync
\item 接收端：Windows + SSH + cwRsync
\item 核心命令：rsync + --rsync-path
\end{itemize}

\end{document}
